\section*{Consensus Search Export}

Manual Consensus searches were first run through the user's logged-in browser session on 2026-06-29. After the Consensus MCP bridge was verified locally, the same assumptions were rechecked through the read-only \texttt{Consensus Academic Research v3.4.0} MCP search tool. The searches were designed to test the specific assumptions behind adopting DynaSpace-style burst-train looming while keeping PPS-kit's 3DTI/SOFA renderer responsible for binaural spatial cues. Browser search pages and MCP-returned paper URLs are listed so the search trail can be re-opened later.

\begin{longtable}{P{0.21\linewidth} P{0.31\linewidth} P{0.23\linewidth} P{0.17\linewidth}}
\caption{Consensus searches and implementation relevance.}\\
\toprule
Question tested & Consensus result & Critical caveat & PPS-kit decision\\
\midrule
\endfirsthead
\toprule
Question tested & Consensus result & Critical caveat & PPS-kit decision\\
\midrule
\endhead
Do approaching looming sounds with rising intensity facilitate faster detection or orienting responses compared with receding or constant sounds?
\newline\url{https://consensus.app/search/approaching-looming-sound-detection/FUKQWKFNSbicvSXmbMchkA/}
& The result was a moderate-to-strong Yes: approaching/rising-intensity sounds are usually more salient and often produce faster orienting or detection responses.
& Some multisensory or visual-dominant tasks show no special looming advantage once task structure and cue type are controlled.
& Keep a strong approach/rising-salience profile, but treat it as a salience-plus-distance manipulation unless matched-level controls are present.\\
\addlinespace
Are repeated broadband noise bursts or onset transients useful for auditory spatial localization compared with continuous narrowband sounds?
\newline\url{https://consensus.app/search/auditory-spatial-localization-cues/85JQWFkfR3-f9kbKCELhBA/}
& The result was Yes: repeated broadband bursts and onset transients usually improve localization relative to continuous narrowband stimuli because they expose richer spectral, binaural, and onset cues.
& Repetition is not always superior to a longer single broadband burst; benefit depends on duration, reverberation, and task.
& Adopt the DynaSpace/Hobeika burst-train source profile as the dry source for the smartphone replication profile.\\
\addlinespace
Are HRTF spectral cues, interaural level differences, and interaural time differences the main cues for binaural spatial hearing and localization?
\newline\url{https://consensus.app/search/binaural-spatial-hearing-cues/3cA7IcmZQD2Vu5c-vJDfLw/}
& Consensus meter: 7 of 7 Yes. The evidence base treats HRTF spectral filtering, ILD, and ITD as the core binaural localization cues, with motion and reverberation adding support.
& Generic HRTFs can be adequate for horizontal lateralization but are weaker for elevation, front/back resolution, and individualized externalization.
& Do not imitate these cues by static spectral hacks in the dry WAV; render them through 3DTI/SOFA and record the chosen HRTF.\\
\addlinespace
Do reverberation, interaural coherence, and direct-to-reverberant ratio affect auditory externalization and distance perception in spatial audio?
\newline\url{https://consensus.app/search/auditory-externalization-spatial-audio/ZupUZceiRGGFHlFGafdDDQ/}
& Consensus meter: 10 of 10 Yes. Reverberation, IACC/coherence structure, and DRR all affect externalization and auditory distance.
& Diotic or poorly structured reverberation can support distance without producing convincing externalization; reverb must be spatially meaningful.
& Add reverb/DRR/IACC as explicit optional QC targets, not as untracked stereo decorrelation.\\
\addlinespace
Are near-field binaural cues important for auditory distance perception and localization for sources within one meter of the head?
\newline\url{https://consensus.app/search/near-field-binaural-auditory-cues/Fz9jdMNQS3e8cbWB0ziT3Q/}
& Consensus meter: 16 papers, 88 percent Yes and 13 percent Mixed. Near-field cues matter under 1 m, especially lateral ILD changes, but are not sufficient alone.
& The strongest near-field effects are lateral/off-midline and interact with level and DRR; ITD changes little with distance.
& Keep the 20 cm endpoint, but prefer native 3DTI near-field-capable rendering for publication stimuli.\\
\addlinespace
Do dynamic approaching sounds facilitate tactile reaction time in peripersonal-space tasks compared with receding or fixed sounds?
\newline\url{https://consensus.app/search/dynamic-sounds-and-tactile-reaction/dw8vS7TcQheii5_C_uOoNQ/}
& The result was Yes with caveats: approaching sounds usually speed tactile responses more than receding or fixed sounds in PPS tasks; prior summaries estimate near-far benefits on the order of tens of milliseconds.
& Expectancy and timing structure can inflate apparent PPS boundaries; Hobeika/Holmes-style critiques require clean controls.
& Use the DynaSpace looming-vs-fixed contrast for smartphone replication, but include matched-level and expectancy-aware variants for stronger inference.\\
\addlinespace
Does individualized or high-quality HRTF rendering improve spatial audio localization and externalization compared with generic non-individual HRTFs?
\newline\url{https://consensus.app/search/individualized-hrtf-spatial-audio/2eIYeqrxTeOiKpEOs8tCNg/}
& The result was mostly Yes: individualized HRTFs improve localization and usually externalization, strongest for elevation, front/back, and VR-like spatial presence.
& Carefully selected generic HRTFs can approach individual performance for some horizontal tasks.
& FABIAN remains an acceptable redistributable default, but the profile should make HRTF identity visible and allow higher-quality/individualized HRTFs later.\\
\addlinespace
Is rising intensity alone sufficient to create auditory looming perception, or do spectral, binaural, and reverberation cues add important spatial information?
\newline\url{https://consensus.app/search/auditory-looming-perception-cues/kEdiFP76QhasleThXrXE-Q/}
& The result was No as a complete account: rising intensity can create classic looming/warning effects, but spectral, binaural, and reverberation cues add spatial information.
& Intensity is often the dominant cue, so a strong gain ramp can masquerade as spatial looming if controls are weak.
& Use rising intensity as one declared component of approach, while preserving independent 3DTI-rendered HRTF, ILD, ITD, near-field, and room cues.\\
\bottomrule
\end{longtable}

\subsection*{Consensus MCP Paper Export}

The following compact export records the read-only MCP searches. The bracket numbers are search-local paper numbers returned by Consensus for that query, not global bibliography numbers. Full abstracts were not copied into this public report; the export preserves query text, result counts, representative titles, and direct Consensus paper URLs.

\begin{longtable}{P{0.20\linewidth} P{0.24\linewidth} P{0.48\linewidth}}
\caption{Consensus MCP query export.}\\
\toprule
Evidence target & Query and count & Representative MCP-returned papers\\
\midrule
\endfirsthead
\toprule
Evidence target & Query and count & Representative MCP-returned papers\\
\midrule
\endhead
Looming salience and warnings
& \texttt{approaching looming sounds rising intensity reaction time detection orienting receding constant sounds}; 20 papers.
& [3] Bach et al. (2009), ``Looming sounds as warning signals: the function of motion cues,'' \url{https://consensus.app/papers/details/45a18fb1f74b5912b30c9465fbb620b3/?utm_source=claude_desktop}. [6] Neuhoff (2001), ``An Adaptive Bias in the Perception of Looming Auditory Motion,'' \url{https://consensus.app/papers/details/0e0c538196b35477ab39922740f634dd/?utm_source=claude_desktop}. [8] Ignatiadis et al. (2024), ``Threat-related corticocortical connectivity elicited by rapid auditory looms,'' \url{https://consensus.app/papers/details/2114c61c357f5e8c854bbc197c06b63a/?utm_source=claude_desktop}.\\
\addlinespace
Broadband bursts and localization
& \texttt{auditory localization broadband noise bursts onset transients narrowband continuous sounds}; 20 papers.
& [1] Middlebrooks (1992), ``Narrow-band sound localization related to external ear acoustics,'' \url{https://consensus.app/papers/details/cef5e85bb5ae529d9543673adde4bf82/?utm_source=claude_desktop}. [2] Kim et al. (2024), ``Median Plane Localization of Short Noise Bursts in the Presence of Horizontally Spread White Noise Masker,'' \url{https://consensus.app/papers/details/dade8599ab9b5fdb849828fd8f9de951/?utm_source=claude_desktop}. [4] Macpherson et al. (2000), ``Localization of brief sounds: effects of level and background noise,'' \url{https://consensus.app/papers/details/b3bdb5d148d550bd8851c3c821390133/?utm_source=claude_desktop}.\\
\addlinespace
Binaural/HRTF cue ownership
& \texttt{binaural spatial hearing HRTF interaural level difference interaural time difference sound localization}; 19 papers.
& [4] Carlini et al. (2024), ``Auditory localization: a comprehensive practical review,'' \url{https://consensus.app/papers/details/ce0b49c63a025ebaa3abe91f3c12f873/?utm_source=claude_desktop}. [3] Gutierrez-Parera et al. (2022), ``Interaural time difference individualization in HRTF by scaling through anthropometric parameters,'' \url{https://consensus.app/papers/details/f79647e2cfc358ef8f3d52a1350d9051/?utm_source=claude_desktop}. [1] Lee et al. (2025), ``Binaural Sound Event Localization and Detection based on HRTF Cues for Humanoid Robots,'' \url{https://consensus.app/papers/details/11d9437d97a45d93b74195c9a49b37c0/?utm_source=claude_desktop}.\\
\addlinespace
Reverb, IACC, DRR, externalization
& \texttt{reverberation interaural coherence direct-to-reverberant ratio auditory externalization distance perception}; 20 papers.
& [1] Choe et al. (2023), ``Impact of direct-to-reverberation ratio and interaural cross-correlation on externalization in binaural audio rendering,'' \url{https://consensus.app/papers/details/e1c815e6c9a55a8db398292d85a3c401/?utm_source=claude_desktop}. [2] Lavandier et al. (2025), ``Reverberation cues underlying virtual auditory distance perception for a frontal source,'' \url{https://consensus.app/papers/details/dd37a407e06a51cbab6251c5ac60e131/?utm_source=claude_desktop}. [5] Kolarik et al. (2013), ``Discrimination of virtual auditory distance using level and direct-to-reverberant ratio cues,'' \url{https://consensus.app/papers/details/9b0b8eedb09d5eff8fbcca667b091eec/?utm_source=claude_desktop}.\\
\addlinespace
Near-field distance cues
& \texttt{near-field binaural cues auditory distance perception within one meter ILD ITD}; 20 papers.
& [2] Brungart et al. (1999), ``Auditory localization of nearby sources. Head-related transfer functions,'' \url{https://consensus.app/papers/details/830a7bfa7d5c5bb4b88e4822debf326d/?utm_source=claude_desktop}. [4] Zhu et al. (2023), ``The contributions of level and near-field head-related transfer functions cues on auditory distance perception in a free field,'' \url{https://consensus.app/papers/details/cd5157c315985b23ab71320b317f770a/?utm_source=claude_desktop}. [5] Guo et al. (2023), ``Distance discrimination thresholds of proximal sound sources in a real anechoic environment,'' \url{https://consensus.app/papers/details/e2c5ce9ed8fd5672aec2443b05c4b141/?utm_source=claude_desktop}.\\
\addlinespace
PPS tactile RT phenomenon
& \texttt{peripersonal space dynamic approaching sounds tactile reaction time fixed receding}; 19 papers.
& [1] Canzoneri et al. (2012), ``Dynamic Sounds Capture the Boundaries of Peripersonal Space Representation in Humans,'' \url{https://consensus.app/papers/details/a947a34ca17b5728b5ae5429d47f8579/?utm_source=claude_desktop}. [2] Matsuda et al. (2021), ``Peripersonal space in the front, rear, left and right directions for audio-tactile multisensory integration,'' \url{https://consensus.app/papers/details/c34b54676bdf50eda323e11dc189b316/?utm_source=claude_desktop}. [3] Hobeika et al. (2019), ``Capturing the dynamics of peripersonal space by integrating sound propagation properties and expectancy effects,'' \url{https://consensus.app/papers/details/53d4e77477fd5a2cad2abd9ea1789dca/?utm_source=claude_desktop}.\\
\addlinespace
Individualized/high-quality HRTF value
& \texttt{individualized HRTF generic HRTF localization externalization spatial audio}; 20 papers.
& [1] Jenny et al. (2020), ``Usability of Individualized Head-Related Transfer Functions in Virtual Reality,'' \url{https://consensus.app/papers/details/f3acff26460c5eadba2bc04664a8a177/?utm_source=claude_desktop}. [2] Frankenbach et al. (2025), ``Evaluating Localization Accuracy of Dynamic Binaural Synthesis,'' \url{https://consensus.app/papers/details/ed8f8111a884542faaac73d52f5f30fa/?utm_source=claude_desktop}. [4] SADIE II case study (2018), ``A Perceptual Evaluation of Individual and Non-Individual HRTFs,'' \url{https://consensus.app/papers/details/c390fc1e332f5147a003b318974b735a/?utm_source=claude_desktop}.\\
\addlinespace
Intensity is not the whole spatial cue set
& \texttt{auditory looming rising intensity spectral cues binaural reverberation distance perception}; 20 papers.
& [3] Kolarik et al. (2015), ``Auditory distance perception in humans: a review of cues,'' \url{https://consensus.app/papers/details/fde5d69ba97851afac0015ca1cf942c9/?utm_source=claude_desktop}. [4] Baumgartner et al. (2017), ``Asymmetries in behavioral and neural responses to spectral cues demonstrate the generality of auditory looming bias,'' \url{https://consensus.app/papers/details/40f6ff7efbcc50f980772b493696877b/?utm_source=claude_desktop}. [5] Mershon et al. (1975), ``Intensity and reverberation as factors in the auditory perception of egocentric distance,'' \url{https://consensus.app/papers/details/af28b552a18a54db9113a715842b4601/?utm_source=claude_desktop}.\\
\bottomrule
\end{longtable}

\subsection*{Consensus-Derived Implementation Claims}

\begin{enumerate}
  \item The dry source should be a burst train. This is supported by the raw DynaSpace files, the Hobeika lineage, and Consensus evidence that broadband transients support localization and attentional salience.
  \item Spatialization should be delegated to 3DTI/SOFA rather than baked into a static waveform. Consensus strongly supports HRTF, ILD, ITD, near-field, reverberation, and DRR as spatial cues that belong in the renderer.
  \item The original smartphone fixed-vs-looming contrast remains a replication target, but it is not a clean causal design by itself. Rising intensity and fixed-source level differences can drive response facilitation unless matched-level control variants are included.
  \item Reverb and low interaural coherence are plausible missing externalization features in the current proxy, but they should be implemented as explicit room/BRIR parameters with QC output.
  \item The toolkit should preserve trajectory flexibility: DynaSpace's 45-degree left approach is one profile, not a hard-coded spatial renderer.
\end{enumerate}
